\section{Proofs for Central Limit Theorems}
\label{sec:annex-p15-central-limit-theorems}

\begin{annexitem}{annex:p15.clt-complex-products}{thm:p15-clt-complex-products}{Product comparison and complex powers}
\begin{proof}
Insert the intermediate products
\(w_1\cdots w_jz_{j+1}\cdots z_n\), \(0\leq j\leq n\).
Successive differences are
\[
 w_1\cdots w_{j-1}(z_j-w_j)z_{j+1}\cdots z_n.
\]
Taking absolute values and summing proves the first inequality,
including \(n=1\), where the empty products are one.
The exponential series is absolutely convergent and, for \(|b|\leq1\),
\[
 |e^b-1-b|\leq |b|^2\sum_{j=2}^\infty\frac1{j!}
 \leq\frac{|b|^2}{2}\sum_{j=0}^\infty2^{-j}=|b|^2.
 \]
Choose \(\gamma>|c|\), increased if necessary to be positive.
Eventually \(|c_n|\leq\gamma\) and \(|c_n|/n\leq1\).
For \(z_j=1+c_n/n\), \(w_j=e^{c_n/n}\), both moduli are at
most \(e^{\gamma/n}\). Hence
\[
 |(1+c_n/n)^n-e^{c_n}|
 \leq e^\gamma n|c_n/n|^2
 \leq e^\gamma\gamma^2/n\to0.
 \]
Continuity of the exponential gives the conclusion. This is the
source's product proof, without replacing it by a logarithm argument.
\end{proof}
\end{annexitem}

\begin{annexitem}{annex:p15.clt-prime-reciprocals}{thm:p15-clt-prime-reciprocals}{An elementary proof of the required prime estimate}
\begin{proof}
We use unique prime factorization only for positive integers.
Its needed elementary justification is as follows. Repeated division
produces a prime factorization by induction on the integer. The
Euclidean algorithm expresses the greatest common divisor of two
integers as an integer linear combination: each remainder is such
a combination, and back-substitution gives this for the last nonzero
remainder. Thus if a prime \(p\) does not divide \(a\), there are
integers \(u,v\) with \(up+va=1\). Multiplication by \(b\)
shows that \(p\mid ab\) implies \(p\mid b\). This prime divisor
property cancels matching primes successively and proves uniqueness.

For every integer \(N\geq1\), each prime \(N<p\leq2N\)
divides \(\binom{2N}{N}\). Indeed it occurs once in \((2N)!\)
and does not occur in either \(N!\). Hence
\[
 \vartheta(2N)-\vartheta(N)
 \leq\log\binom{2N}{N}\leq N\log4,
 \]
where the last bound follows because \(\binom{2N}{N}\)
is one nonnegative summand of
\(\sum_{j=0}^{2N}\binom{2N}{j}=2^{2N}\).
Summing with \(N=1,2,\ldots,2^{k-1}\) through powers of two
gives \(\vartheta(2^k)\leq(2^k-1)\log4\).
For \(x\geq2\), choose \(2^{k-1}<x\leq2^k\); then
\(\vartheta(x)\leq2x\log4\).
Put \(C=2\log4\). Summing
\(1/p=1/x+\int_p^x u^{-2}\,du\), after multiplying by
\(\log p\), yields the exact identity
\[
 \sum_{p\leq x}\frac{\log p}{p}
 =\frac{\vartheta(x)}x+\int_2^x\frac{\vartheta(u)}{u^2}\,du
 \leq C(1+\log x).
 \]
All sums here are finite, so the interchange needs no convergence
assumption.

Let \(A(x)=\sum_{p\leq x}1/p\). The finite Euler product
\[
 P(x)=\prod_{p\leq x}(1-1/p)^{-1}
 \]
expands by nonnegative geometric series into the reciprocals
of all integers having no prime factor above \(x\), each exactly
once. In particular \(P(x)\geq\sum_{j\leq\lfloor x\rfloor}1/j\).
Also, by the real logarithm series,
\[
 0\leq\log P(x)-A(x)
 =\sum_{p\leq x}\sum_{r\geq2}\frac1{rp^r}
 \leq2\sum_{j\geq2}j^{-2}<\infty
 \]
with a bound independent of \(x\). Integral comparison of
the harmonic sum with \(\log x\) proves
\(A(x)\geq\log\log x-O(1)\).

For the reverse bound take \(\delta=1/\log x\) and \(s=1+\delta\).
For \(p\leq x\),
\[
 0\leq p^{-1}-p^{-s}
 =p^{-1}(1-e^{-\delta\log p})
 \leq\delta\log p/p.
 \]
The preceding bound makes the sum of these differences \(O(1)\).
The finite Euler product at exponent \(s>1\) is at most
\(\sum_{j\geq1}j^{-s}\), by the same nonnegative expansion.
Since \(u^{-s}\) decreases,
\[
 \sum_{j\geq1}j^{-s}\leq1+\int_1^\infty u^{-s}\,du
                           =1+\delta^{-1}=1+\log x.
 \]
Taking logarithms of the product gives
\[
 \sum_{p\leq x}p^{-s}
 \leq\log\prod_{p\leq x}(1-p^{-s})^{-1}
 \leq\log(1+\log x)=\log\log x+O(1).
 \]
Together the estimates give \(A(x)=\log\log x+O(1)\).
Finally \(\sum_pp^{-2}\leq\sum_{j\geq2}j^{-2}<\infty\).
This proves all number-theoretic premises actually used below.
\end{proof}
\end{annexitem}

\begin{annexitem}{annex:p15.clt-erdos-kac}{thm:p15-clt-erdos-kac}{The source's cutoff and complete moment transfer}
\begin{proof}
For each prime \(p\), the proportion of \(1\leq m\leq n\)
divisible by \(p\) is \(\lfloor n/p\rfloor/n\to1/p\).
For finitely many distinct primes it is
\(\lfloor n/(p_1\cdots p_k)\rfloor/n\), tending to
\((p_1\cdots p_k)^{-1}\). These limits explain the independent
comparison but do not assert independence under \(\mathbb P_n\).
They also do not define a countably additive law on
\(\mathbb N\): each singleton has limiting mass zero whereas
the entire space has limiting mass one.

For sufficiently large \(n\), put
\[
 \ell_n=\log\log n,\qquad
 \alpha_n=n^{1/\ell_n},\qquad
 g_n(m)=\sum_{p\leq\alpha_n}\delta_p(m).
 \]
The cutoff has
\(\log\log\alpha_n=\ell_n-\log\ell_n\), and
\(\alpha_n^r/n\to0\) for every fixed finite \(r>0\), because
its logarithm is \((r/\ell_n-1)\log n\to-\infty\).
The prime reciprocal estimate gives
\[
 \sum_{\alpha_n<p\leq n}\frac1p
      =\log\ell_n+O(1)=o(\sqrt{\ell_n}).
 \]
Thus the nonnegative difference \(g-g_n\), under
\(\mathbb P_n\), has expectation at most this sum. Markov's
inequality proves
\((g-g_n)/\sqrt{\ell_n}\to0\) in probability.

Construct independent Bernoulli variables \(B_p\) of probabilities
\(1/p\) on the supplied countable product space, and set
\[
 V_n=\sum_{p\leq\alpha_n}B_p,\quad
 b_n=\mathbb EV_n=\sum_{p\leq\alpha_n}1/p,\quad
 a_n^2=\operatorname{Var}(V_n)
      =\sum_{p\leq\alpha_n}(1/p)(1-1/p).
 \]
Then
\[
 b_n=\ell_n-\log\ell_n+O(1),\qquad
 a_n^2=\ell_n-\log\ell_n+O(1).
 \]
In particular \(a_n\to\infty\),
\(a_n/\sqrt{\ell_n}\to1\), and
\((b_n-\ell_n)/\sqrt{\ell_n}\to0\).
The independent centered rows \((B_p-1/p)/a_n\) have total
variance one and uniform absolute bound \(1/a_n\to0\).
The full moment-array theorem gives
\[
 \mathbb E\{(V_n-b_n)/a_n\}^r\to\mathbb E\chi^r
                                      \quad(r\geq1\text{ integer}).
 \]

We next transfer every one of these moments to \(g_n\).
Expand the \(r\)-th power as a sum over ordered \(r\)-tuples
of primes at most \(\alpha_n\), allowing repetitions.
For such a tuple let \(q\) be the product of its distinct primes.
Bernoulli idempotence and independence give expectation \(1/q\);
the corresponding divisibility product has expectation
\(\lfloor n/q\rfloor/n\). Their difference in absolute value
is at most \(1/n\), whether or not \(q>n\).
There are at most \(\alpha_n^r\) tuples, so
\[
 |\mathbb EV_n^r-\mathbb E_ng_n^r|
                         \leq\alpha_n^r/n.
 \]
At \(r=0\) the two moments are exactly one.
The binomial expansion of the centered \(r\)-th powers consequently
gives
\[
 |\mathbb E(V_n-b_n)^r-\mathbb E_n(g_n-b_n)^r|
 \leq \frac1n\sum_{j=0}^r\binom rj
                  \alpha_n^j b_n^{r-j}
 =\frac{(\alpha_n+b_n)^r}{n}\longrightarrow0.
 \]
Indeed \(0\leq b_n\leq\#\{p\leq\alpha_n\}\leq\alpha_n\),
and \(\alpha_n^r/n\to0\). Division by \(a_n^r\), which
is at least one eventually, transfers the normalized moments.
The normal law is moment-determinate, so the method of moments
implies \((g_n-b_n)/a_n\Rightarrow\chi\) under \(\mathbb P_n\).
The deterministic changes of center and scale, and the negligible
discarded prime sum, give the stated limit for \(g\).
This proves the source's complete moment-comparison route.
\end{proof}
\end{annexitem}

\begin{annexitem}{annex:p15.clt-smoothing-bound}{thm:p15-clt-smoothing-bound}{The source's symmetric smoothing estimate}
\begin{proof}
The Pólya density and transform were proved in the characteristic
chapter; scaling a variable with that density by \(1/L\)
gives \(h_L\) and \(\omega_L\). It is even, integrates to one,
and satisfies \(h_L(x)\leq2/(\pi Lx^2)\) for \(x\ne0\).
The value at zero follows from the cosine expansion.
Product integration gives the displayed convolution CDF identity.

If \(\eta=0\), the bound is immediate. Suppose \(\eta>0\).
It suffices to consider the case that the positive supremum of
\(D=F-G\) is \(\eta\).
To justify that reduction, reflecting the laws gives CDFs
\(\widetilde F(x)=1-F(-x-)\),
\(\widetilde G(x)=1-G(-x)\), so their difference is
\(-D(-x-)\). Right continuity and the existence of left limits
show that the supremum of absolute left-limit values equals
the original supremum. Reflection therefore interchanges the
positive and negative supremum. Since \(H_L\) has a density
and is symmetric, both convolved CDFs are continuous, and their
reflected difference is \(-(D*H_L)(-x)\). Its supremum norm
is unchanged.

For arbitrary \(e>0\), choose \(x_0\) with
\(D(x_0)\geq\eta-e\), avoiding any assumption that the supremum
is attained. Nondecrease of \(F\) and the Lipschitz bound for \(G\)
give \(D(x_0+s)\geq\eta-e-\lambda s\) for \(s\geq0\).
Put \(\delta=\eta/(2\lambda)>0\) and \(t=x_0+\delta\).
For \(|x|\leq\delta\), this says
\[
 D(t-x)\geq\eta/2-e+\lambda x.
 \]
Everywhere \(D\geq-\eta\). The tail mass \(q\) of \(h_L\)
outside \([-\delta,\delta]\) obeys
\[
 q\leq2\int_\delta^\infty\frac{2}{\pi Lx^2}\,dx
                          =\frac4{\pi L\delta}.
 \]
Symmetry gives \(\int_{-\delta}^{\delta}xh_L(x)\,dx=0\).
Consequently
\[
 (D*H_L)(t)
 \geq(\eta/2-e)(1-q)-\eta q
 \geq\eta/2-e-\frac{6\eta}{\pi L\delta}
 =\eta/2-e-\frac{12\lambda}{\pi L}.
 \]
The estimate remains valid if the displayed upper bound on \(q\)
exceeds one. Let \(e\downarrow0\), then rearrange.
If \(G\) is differentiable with derivative at most \(\lambda\),
the mean value theorem gives exactly the Lipschitz premise.
\end{proof}
\end{annexitem}

\begin{annexitem}{annex:p15.clt-cdf-inversion}{thm:p15-clt-cdf-inversion}{The integrable quotient, including the convolved laws}
\begin{proof}
Integrable characteristic functions have continuous densities
\[
 k_i(y)=\frac1{2\pi}\int_{\mathbb R}e^{-ity}\kappa_i(t)\,dt
 \]
by the proved integrable inversion theorem.
For finite \(a<x\), absolute integrability on the finite
interval in \(y\) justifies Fubini and gives
\[
 (K_1-K_2)(x)-(K_1-K_2)(a)
 =\frac1{2\pi}\int_{\mathbb R}
   (e^{-ita}-e^{-itx})\frac{\kappa_1(t)-\kappa_2(t)}{it}\,dt.
 \]
The quotient is integrable by hypothesis. Its oscillatory integral
with \(e^{-ita}\) tends to zero as \(a\to-\infty\) by the
proved Riemann--Lebesgue lemma. Both CDFs tend to zero there,
so the claimed formula follows.

If both laws have finite first absolute moments, then
\[
 |\kappa_i(t)-1|\leq |t|\,\mathbb E|Z_i|
 \]
for a variable \(Z_i\) with that law. The difference quotient
is bounded near zero. For \(|t|\geq1\), it is bounded by
\(|\kappa_1(t)|+|\kappa_2(t)|\), an integrable function.
This proves the stated sufficient condition, including the
source's mean-zero case. In that case differentiability even
gives the quotient tending to zero at zero.

For the smoothing application take
\(\kappa_1=\varphi_F\omega_L\),
\(\kappa_2=\varphi_G\omega_L\).
They are bounded and supported on \([-L,L]\), hence integrable.
Their difference divided by \(t\) is integrable: near zero
the preceding first-moment bound applies to the original laws
\(F,G\), and away from zero bounded compact support suffices.
We have not assumed finite first moments of \(F*H_L\) or
\(G*H_L\). The inversion identity thus gives
\[
 \eta_L\leq\frac1{2\pi}\int_{-L}^L
          \frac{|\varphi_F(t)-\varphi_G(t)|}{|t|}\,dt,
 \]
since \(0\leq\omega_L\leq1\).
Combine this with the preceding smoothing lemma. This validates
the entire source application despite the infinite absolute first
moment of the smoothing distribution itself. To see that this
moment is indeed infinite, substitute \(u=Lx\) in
\(\int_0^\infty xh_L(x)\,dx\). Apart from a positive constant
the integral is \(\int_0^\infty(1-\cos u)/u\,du\).
On each interval
\([2k\pi+\pi/2,2k\pi+3\pi/2]\), \(k\geq1\), its numerator
is at least one, so the contributions dominate a positive
constant times \(\sum_{k\geq1}1/k=\infty\).
If \(Y\) has this smoothing law and \(X\) is integrable,
\(|Y|\leq|X+Y|+|X|\) forces \(\mathbb E|X+Y|=\infty\).
\end{proof}
\end{annexitem}

\begin{annexitem}{annex:p15.clt-berry-esseen}{thm:p15-clt-berry-esseen}{The full cutoff and constant calculation}
\begin{proof}
Scale \(X_j\) by \(\sigma\). It is enough to prove the result
at variance one, writing \(\rho\) for the scaled third absolute
moment. Hölder's inequality gives
\(1=\mathbb EX_1^2\leq(\mathbb E|X_1|^3)^{2/3}\), so
\(\rho\geq1\). For \(n\leq9\), \(3\rho/\sqrt n\geq1\),
and a CDF distance is at most one. Hence assume \(n\geq10\).

Let \(\varphi(t)=\mathbb Ee^{itX_1}\) and
\(L=4\sqrt n/(3\rho)\).
The proved Taylor bound with its factorial constant gives
\[
 |\varphi(t)-1+t^2/2|\leq\rho|t|^3/6.
 \]
For \(|\theta|\leq L\), we have
\(|\theta|/\sqrt n\leq4/(3\rho)\leq4/3<\sqrt2\).
Therefore \(1-\theta^2/(2n)\geq0\) and
\[
 |\varphi(\theta/\sqrt n)|
 \leq1-\frac{\theta^2}{2n}
                  +\frac{\rho|\theta|^3}{6n^{3/2}}
 \leq1-\frac{5\theta^2}{18n}
 \leq e^{-5\theta^2/(18n)}.
 \]
Set
\(\alpha=\varphi(\theta/\sqrt n)\),
\(\beta=e^{-\theta^2/(2n)}\), and
\(\gamma=e^{-5\theta^2/(18n)}\).
Both moduli are bounded by \(\gamma\). The product comparison
gives \(|\alpha^n-\beta^n|\leq n|\alpha-\beta|\gamma^{n-1}\),
and \(n\geq10\) gives
\(\gamma^{n-1}\leq e^{-\theta^2/4}\).

For \(0\leq x\leq1\), the alternating exponential series has
decreasing absolute terms, so the source's alternating-series
estimate gives \(|e^{-x}-1+x|\leq x^2/2\).
This applies here because
\(\theta^2/(2n)\leq8/(9\rho^2)<1\).
The bound also holds for every \(x\geq0\), by integrating
the second derivative \(e^{-x}\leq1\) in its Taylor remainder.
Consequently
\[
 n|\alpha-\beta|
 \leq\frac{\rho|\theta|^3}{6\sqrt n}+\frac{\theta^4}{8n}.
 \]
For \(\theta\ne0\), combine the last two estimates:
\[
 \frac{|\varphi(\theta/\sqrt n)^n-e^{-\theta^2/2}|}{|\theta|}
 \leq e^{-\theta^2/4}
       \left(\frac{\rho\theta^2}{6\sqrt n}
                         +\frac{|\theta|^3}{8n}\right)
 \leq\frac1L e^{-\theta^2/4}
                      \left(\frac{2\theta^2}{9}
                                      +\frac{|\theta|^3}{18}\right).
 \]
For the last inequality use \(\rho/\sqrt n=4/(3L)\) and
\(1/n\leq4/(9L)\), the latter following from
\(L\leq4\sqrt n/3\) and \(\sqrt n\geq3\).
At zero all terms have continuous value zero.

The normal CDF is Lipschitz with
\(\lambda=(2\pi)^{-1/2}<2/5\).
Apply the smoothing inequality to \(F_n,\Phi\).
All first absolute moments are finite, so its premises hold.
Enlarge the nonnegative integral from \([-L,L]\) to
\(\mathbb R\) to get
\[
 \pi L\,d(F_n,\Phi)
 \leq
 \int_{\mathbb R} e^{-\theta^2/4}
       \left(\frac{2\theta^2}{9}+\frac{|\theta|^3}{18}\right)\,d\theta
                                           +\frac{48}{5}.
 \]
The already proved normal normalization and variance, at variance
two, give \(\int_{\mathbb R}\theta^2e^{-\theta^2/4}\,d\theta
=4\sqrt\pi\).
The substitution \(u=\theta^2/4\) on the positive half-line gives
\[
 \int_{\mathbb R}|\theta|^3e^{-\theta^2/4}\,d\theta
 =16\int_0^\infty ue^{-u}\,du=16,
 \]
where integration by parts makes the last integral one.
The source's direct integration-by-parts route gives the same
value: writing \(x^3=x^2x\) and
\(xe^{-x^2/4}\,dx=-2\,d(e^{-x^2/4})\),
\[
 2\int_0^\infty x^3e^{-x^2/4}\,dx
 =-4[x^2e^{-x^2/4}]_0^\infty
                         +8\int_0^\infty xe^{-x^2/4}\,dx=16.
 \]
It follows that
\[
 d(F_n,\Phi)
 \leq \frac{3}{4\pi}
       \left(\frac{8\sqrt\pi}{9}+\frac89+\frac{48}{5}\right)
       \frac{\rho}{\sqrt n}
 <3\,\frac{\rho}{\sqrt n}.
 \]
Here is an exact check of the final strict inequality without a
decimal calculation. It is equivalent to
\(8\sqrt\pi/9+472/45<4\pi\).
For \(x\geq25/8\), the function
\(4x-8\sqrt{x}/9-472/45\) has positive derivative
\(4-4/(9\sqrt x)>0\).
At \(x=25/8\), use \(\sqrt{25/8}<9/5\); its value is greater than
\[
 \frac{25}{2}-\frac85-\frac{472}{45}
 =\frac{37}{90}>0.
 \]
The elementary bound \(\pi>25/8\) follows, for example, from
the perimeter of a regular inscribed \(24\)-gon in the unit
circle: \(\pi>24\sin(\pi/24)>25/8\).
To check the last algebraic bound, half-angle identities give
\(\sin^2(\pi/24)=(4-\sqrt6-\sqrt2)/8\);
using \(\sqrt6<24495/10000\) and
\(\sqrt2<14143/10000\), verified by squaring, makes this
greater than \((25/192)^2\).
The same \(\pi>25/8\) also proves
\((2\pi)^{-1/2}<2/5\), used above.
Rescaling returns the factor \(\rho/\sigma^3\).

Finally take fair-sign summands and an even sample size \(2m\).
Symmetry gives
\[
 F_{2m}(0)-\Phi(0)
 =\tfrac12\mathbb P(S_{2m}=0)
 =\tfrac12\binom{2m}{m}2^{-2m}
 \sim\frac1{2\sqrt{\pi m}},
 \]
by the already proved Stirling formula. This lower bound at
one point proves the claimed obstruction to a uniformly smaller
order than \(n^{-1/2}\), even for this bounded law.
\end{proof}
\end{annexitem}

\begin{annexitem}{annex:p15.clt-iid}{thm:p15-clt-iid}{The characteristic-function proof}
\begin{proof}
Replace \(X_j\) by \(X_j-\mu\). For the centered common law,
the proved moment expansion
(Theorem~\ref{thm:p15-characteristic-moment-expansion}) gives
\[
 \varphi(t)=1-\sigma^2t^2/2+o(t^2).
 \]
Independence and the affine transform identity give
\[
 \mathbb Ee^{itS_n/(\sigma\sqrt n)}
 =\varphi(t/(\sigma\sqrt n))^n
 =\{1-t^2/(2n)+o(n^{-1})\}^n
 \longrightarrow e^{-t^2/2}.
 \]
The complex-power result just proved applies for each fixed real \(t\).
The limit is the characteristic function of the already constructed
standard normal law. The continuity theorem therefore proves weak
convergence. If \(\sigma=0\), the nonnegative variable
\((X_j-\mu)^2\) has expectation zero and is zero almost surely.
A countable intersection gives all centered sums zero on one
event of probability one.
\end{proof}
\end{annexitem}

\begin{annexitem}{annex:p15.clt-perturbation}{thm:p15-clt-perturbation}{Tightness and the distribution-function sandwich}
\begin{proof}
At every real \(x\) and every \(\varepsilon>0\),
\[
 \mathbb P(Z_n\leq x-\varepsilon)-\mathbb P(|R_n|>\varepsilon)
 \leq\mathbb P(Z_n+R_n\leq x)
 \leq\mathbb P(Z_n\leq x+\varepsilon)
                         +\mathbb P(|R_n|>\varepsilon).
 \]
At a continuity point of the CDF of \(Z\), choose
\(\varepsilon\downarrow0\) so that both endpoints are continuity
points; this is possible because a CDF has at most countably many
discontinuities. First let \(n\to\infty\), then
\(\varepsilon\downarrow0\). This proves the perturbation assertion.

Weak convergence makes the laws of \(Z_n\) tight. Explicitly, for
arbitrarily large continuity points \(K,-K\) of the limiting CDF,
its tail probabilities are small; weak convergence bounds the
corresponding tails for large \(n\), and enlarging \(K\) handles
the finitely many earlier laws. Thus
\[
 \mathbb P(|(A_n-a)Z_n|>\varepsilon)
 \leq\mathbb P(|Z_n|>K)+\mathbb P(|A_n-a|>\varepsilon/K)\to0
 \]
after taking first \(n\to\infty\), then \(K\to\infty\).
The deterministic continuous mapping \(z\mapsto az+b\) preserves
weak convergence, by the bounded-continuous characterization.
Apply the first assertion to
\((A_n-a)Z_n+(B_n-b)\).
If \(a\ne0\), then
\(\mathbb P(|A_n|\leq|a|/2)\to0\), and
\[
 |A_n^{-1}-a^{-1}|
 \leq2|A_n-a|/|a|^2\quad\text{when }|A_n|\geq|a|/2.
 \]
This proves the reciprocal version and makes the assigned values
on \(\{A_n=0\}\) immaterial.
Finally, a constant \(d_n\) has characteristic function of modulus
one. Pointwise convergence of those functions to
\(e^{itb-\sigma^2t^2/2}\), with \(\sigma>0\), is impossible
at any \(t\ne0\).
\end{proof}
\end{annexitem}

\begin{annexitem}{annex:p15.clt-numerical-examples}{thm:p15-clt-numerical-examples}{The finite calculations and the real Poisson parameter}
\begin{proof}
For roulette,
\(\mathbb EX=(18-20)/38=-1/19\) and
\(\mathbb EX^2=1\). Hence
\(\sigma^2=1-19^{-2}=360/361\).
The standardized threshold for \(n=361\) is
\[
 \frac{-n\mu}{\sigma\sqrt n}
 =\frac{19}{\sqrt{360}}.
 \]
The normal approximation substitutes its upper tail for the
corresponding standardized sum tail. At odd \(361\), zero itself
is not a possible sum; this does not turn the approximation into
an exact lattice probability. The expected gain is \(361\mu=-19\).

For a Bernoulli variable of parameter \(p\),
\(\mathbb EX=p\) and
\(\operatorname{Var}(X)=p(1-p)\), by \(X^2=X\).
At \(p=1/2\) these are \(1/2,1/4\). Thus the standardized
endpoints for \(4900\) and \(5100\) are \(-2,2\).
At \(n=16\), the correction interval for eight successes is
\([7.5,8.5]\), which standardizes to \([-1/4,1/4]\).
Independence gives the exact mass
\(\binom{16}{8}2^{-16}=12870/65536\).
For the die calculation, the six-indicators have parameter \(1/6\),
so the mean is \(180/6=30\) and the variance is
\(180(1/6)(5/6)=25\).
Fewer than \(25\) means at most \(24\), giving the corrected
upper threshold \(24.5\) and standardized threshold \(-1.1\).
The stated decimals are rounded evaluations of the displayed
normal integrals. None is claimed to follow as an exact
finite-sample equality from an asymptotic theorem.

For completeness the Poisson mass formula
\(e^{-\lambda}\lambda^j/j!\), \(j\geq0\), gives
\(\mathbb EZ_\lambda=\lambda\) and
\(\mathbb E[Z_\lambda(Z_\lambda-1)]=\lambda^2\) by shifting
the absolutely convergent series. Its variance is therefore
\(\lambda\). If independent Poisson variables have means
\(a,b\geq0\), their convolution at \(j\) is
\[
 e^{-a-b}\sum_{\ell=0}^j
       \frac{a^\ell b^{j-\ell}}{\ell!(j-\ell)!}
 =e^{-a-b}\frac{(a+b)^j}{j!};
 \]
at mean zero the law is point mass at zero.
Consequently a sum of \(n\) independent mean-one Poisson
variables is Poisson of mean \(n\), and the i.i.d. CLT applies.

For general \(\lambda\to\infty\), set \(n=\lfloor\lambda\rfloor\).
Construct independent Poisson variables of means \(n\),
\(\lambda-n\), and \(n+1-\lambda\). Their successive sums
couple \(Z_n,Z_\lambda,Z_{n+1}\) with
\(Z_n\leq Z_\lambda\leq Z_{n+1}\).
Both
\[
 (Z_n-\lambda)/\sqrt\lambda
 =\sqrt{n/\lambda}\,(Z_n-n)/\sqrt n+(n-\lambda)/\sqrt\lambda
 \]
and the analogous expression with \(n+1\) converge to \(\chi\).
The CDF inequalities between the three coupled variables give
the same limit for \(Z_\lambda\), since \(\Phi\) is continuous.
This retains the source's monotone coupling argument.
\end{proof}
\end{annexitem}

\begin{annexitem}{annex:p15.clt-pairwise-boundary}{thm:p15-clt-pairwise-boundary}{The sign-product construction}
\begin{proof}
Take independent fair signs \(\xi_1,\xi_2,\ldots\).
Set \(X_1=\xi_1\). Recursively, once
\(X_1,\ldots,X_{2^r}\) have been defined, put
\[
 X_{2^r+j}=X_j\xi_{r+2}\quad(1\leq j\leq2^r),\qquad r\geq0.
 \]
The first \(2^r\) variables are precisely
\(\xi_1\prod_{\ell\in A}\xi_\ell\) as \(A\) runs over the
subsets of \(\{2,\ldots,r+1\}\). Induction proves this statement,
and hence the recursion is well-defined for all positive indices.
Each is a fair sign. For distinct such monomials \(X_j,X_k\),
their product is a nonempty product of independent fair signs,
so its expectation is zero. For \(a,b\in\{-1,1\}\),
\[
 \mathbb P(X_j=a,X_k=b)
 =\mathbb E[(1+aX_j)(1+bX_k)]/4=1/4.
 \]
This proves pairwise independence directly.
Expansion of a finite product gives
\[
 S_{2^r}=\xi_1\prod_{\ell=2}^{r+1}(1+\xi_\ell).
 \]
It is nonzero only when all the latter \(r\) signs are \(1\),
an event of probability \(2^{-r}\), independently of \(\xi_1\).
The stated masses follow. The probability that its normalization
is nonzero tends to zero, proving convergence in probability.
All \(X_j\) have variance one, so pairwise independence still
gives \(\operatorname{Var}(S_n)=n\).
\end{proof}
\end{annexitem}

\begin{annexitem}{annex:p15.clt-oscillation}{thm:p15-clt-oscillation}{Tail events and the factorial subsequence}
\begin{proof}
Let \(L=\limsup_n S_n/\sqrt n\), an extended measurable variable.
For every fixed \(k\), the finite prefix \(S_k/\sqrt n\to0\)
pointwise, so
\[
 L=\limsup_n\frac{X_{k+1}+\cdots+X_n}{\sqrt n}.
 \]
Thus \(\{L\geq b\}\) belongs to every tail sigma-algebra.
For fixed \(b\in\mathbb R\), the event
\(\{S_n/\sqrt n>b\text{ infinitely often}\}\) implies \(L\geq b\).
Its probability is at least
\(\limsup_n\mathbb P(S_n/\sqrt n>b)=1-\Phi(b/\sigma)>0\):
each union over \(n\geq N\) has probability at least the supremum
of its individual probabilities, then let \(N\to\infty\).
Kolmogorov's tail zero-one law makes
\(\mathbb P(L\geq b)=1\). Intersect over positive integer \(b\).
Applying the same argument to \(-X_j\) gives the liminf assertion.

For the stronger failure of convergence in probability, suppose
\(S_n/\sqrt n\to Z\) in probability. Then \(Z\) has law
\(N(0,\sigma^2)\), by the CLT and uniqueness of a weak limit.
For \(m\geq2\), put
\[
 A_m=\frac{S_{(m-1)!}}{\sqrt{(m-1)!}},\qquad
 B_m=\frac{S_{m!}-S_{(m-1)!}}{\sqrt{m!}}.
 \]
The two variables are independent. Also \(A_m\to Z\) in
probability, and
\[
 B_m=S_{m!}/\sqrt{m!}-A_m/\sqrt m\longrightarrow Z
                                               \quad\text{in probability},
 \]
because \(A_m\) is tight. Hence \(B_m-A_m\to0\) in probability.
For fixed \(t\), boundedness and convergence in probability imply
convergence of the expectations of the exponentials: split the
expectation on a small deviation event and its complement, using
\(|e^{iu}-e^{iv}|\leq\min(2,|u-v|)\).
Independence would therefore imply simultaneously
\[
 \mathbb Ee^{it(B_m-A_m)}\to1,\qquad
 \mathbb Ee^{it(B_m-A_m)}
 =\mathbb Ee^{itB_m}\mathbb Ee^{-itA_m}
 \longrightarrow e^{-\sigma^2t^2}.
 \]
These contradict each other for \(t\ne0\).
\end{proof}
\end{annexitem}

\begin{annexitem}{annex:p15.clt-sqrt-scale-necessity}{thm:p15-clt-sqrt-scale-necessity}{Independent-copy symmetrization and the conditional signs}
\begin{proof}
Suppose first that \(\mathbb EX_1^2=\infty\).
On a product space take an independent copy \(X'_j\) and put
\(Y_j=X_j-X'_j\). These are i.i.d. symmetric variables.
Choose \(M<\infty\) with \(p=\mathbb P(|X'_1|\leq M)>0\).
Independence gives
\[
 \mathbb EY_1^2\geq
 \frac p4\,\mathbb E(X_1^2;|X_1|>2M)=\infty.
 \]
Let \(U_j=Y_j\mathbf1_{\{|Y_j|\leq A\}}\) and
\(V_j=Y_j\mathbf1_{\{|Y_j|>A\}}\).
The \(U_j\) are centered, bounded, i.i.d.; their variances
\(\tau_A^2=\mathbb E(Y_1^2;|Y_1|\leq A)\) increase to infinity.

Here \(U_j\) and \(V_j\) are not asserted independent.
To justify the needed inequality, a symmetric law can be
represented as \(\varepsilon R\), with \(R\) of law \(|Y_1|\)
and \(\varepsilon\) an independent fair sign. Indeed the two
reflected half-laws give the original symmetric law, including
any atom at zero. Use independent such pairs for the whole sequence.
For each fixed \(n\), condition by integration over the magnitudes
\(R_1,\ldots,R_n\) and the signs with \(R_j\leq A\).
The remaining signs are still independent and fair, so their
weighted sum \(\sum V_j\) is symmetric and has probability at
least \(1/2\) of being nonnegative. Integrating this finite product
calculation gives, for \(K>0\),
\[
 \mathbb P\left(\sum_{j=1}^nY_j\geq K\sqrt n\right)
 \geq\frac12\mathbb P\left(\sum_{j=1}^nU_j\geq K\sqrt n\right).
 \]
Choose \(A\) so large that
\(1-\Phi(K/\tau_A)>2/5\).
The bounded i.i.d. CLT makes the last right side greater than
\(1/5\) for all sufficiently large \(n\), with a slightly larger
choice of \(A\) if necessary to make the inequality strict.

But weak convergence of \(S_n/\sqrt n\) makes that family tight.
Its independent-copy differences are tight too, since
\[
 \mathbb P(|S_n-S'_n|>2K\sqrt n)
 \leq2\mathbb P(|S_n|>K\sqrt n).
 \]
The preceding lower bound for every \(K\) contradicts this.
Thus \(\mathbb EX_1^2<\infty\).
Its mean \(\mu\) exists, and the weak law gives \(S_n/n\to\mu\)
in probability. Tightness of \(S_n/\sqrt n\) also gives
\(S_n/n\to0\) in probability. Hence \(\mu=0\).
The i.i.d. CLT, including its zero-variance case, identifies the
only possible weak limit.
\end{proof}
\end{annexitem}

\begin{annexitem}{annex:p15.clt-transformed-sums}{thm:p15-clt-transformed-sums}{The two random-denominator calculations}
\begin{proof}
For the nonnegative observations, the weak law gives \(S_n/n\to1\)
in probability. Rationalization gives the everywhere defined identity
\[
 2(\sqrt{S_n}-\sqrt n)
 =\frac{S_n-n}{\sqrt n}\,
                      \frac{2}{\sqrt{S_n/n}+1}.
 \]
The first factor converges weakly to \(\sigma\chi\) and the
second tends in probability to one. The perturbation lemma applies.

For the self-normalized statistic, the integrable weak law applied
to \(X_j^2\) gives \(n^{-1}\sum_jX_j^2\to\sigma^2\) in
probability. Its nonnegative square root tends to \(\sigma>0\).
Divide the CLT for \(S_n/\sqrt n\) by this root. In particular
the probability that the denominator vanishes tends to zero,
so the specified finite value there has no effect on weak convergence.
\end{proof}
\end{annexitem}

\begin{annexitem}{annex:p15.clt-random-index}{thm:p15-clt-random-index}{Two maximal inequalities around a deterministic index}
\begin{proof}
Fix \(0<\delta<1/2\) and \(\varepsilon>0\). Write \(a=a_n\)
and \(r=\lceil\delta a\rceil\); for sufficiently large \(n\),
\(r\leq a\). On \(|N_n/a-1|\leq\delta\), the difference
\(|S_{N_n}-S_a|\) is bounded by the larger of the maxima
\[
 \max_{1\leq j\leq r}|X_{a+1}+\cdots+X_{a+j}|,\qquad
 \max_{1\leq j\leq r}|X_a+\cdots+X_{a-j+1}|.
 \]
Each displayed sequence of increments is independent and centered
with variance \(\sigma^2\) per term, even in reversed order.
Kolmogorov's maximal inequality gives
\[
 \mathbb P(|S_{N_n}-S_a|>\varepsilon\sigma\sqrt a)
 \leq\mathbb P(|N_n/a-1|>\delta)
                                  +\frac{2r}{\varepsilon^2a}.
 \]
Take the limsup in \(n\), then let \(\delta\downarrow0\).
The normalized difference tends to zero in probability.
The CLT for \(S_a/(\sigma\sqrt a)\), followed by perturbation,
proves the theorem. Every estimate depends only on the displayed
index and convergences, so it holds equally along every real
index tending to infinity. This is also the sequential criterion
for the stated continuously indexed limit.
\end{proof}
\end{annexitem}

\begin{annexitem}{annex:p15.clt-renewal-random-index}{thm:p15-clt-renewal-random-index}{Inverting the centered sum with an overshoot bound}
\begin{proof}
The renewal inverse law gives \(N(t)/t\to1/\mu\) almost surely.
For \(a_t=\lfloor t/\mu\rfloor\), eventually positive,
\(N(t)/a_t\to1\). Apply the random-index theorem to the centered
increments \(Y_j-\mu\):
\[
 \frac{T_{N(t)}-\mu N(t)}{\sigma\sqrt{t/\mu}}
                                      \Rightarrow\chi.
 \]
We must control the difference between \(T_{N(t)}\) and \(t\).
For any fixed \(C>1/\mu\),
\[
 \mathbb P(Y_{N(t)+1}>\varepsilon\sqrt t)
 \leq\mathbb P(N(t)>Ct)
       +(\lfloor Ct\rfloor+1)\mathbb P(Y_1>\varepsilon\sqrt t).
 \]
The first term tends to zero. Finite second moment implies
\[
 t\,\mathbb P(Y_1>\varepsilon\sqrt t)
 \leq\varepsilon^{-2}
       \mathbb E(Y_1^2;Y_1>\varepsilon\sqrt t)\longrightarrow0
 \]
by dominated convergence. On the event defining the finite count,
\(0\leq t-T_{N(t)}<Y_{N(t)+1}\), including zero waiting times.
The overshoot is thus \(o_{\mathbb P}(\sqrt t)\). Finally
\[
 \mu N(t)-t
 =-(T_{N(t)}-\mu N(t))-(t-T_{N(t)}).
 \]
Perturbation and the symmetry of \(\chi\) give the result.
\end{proof}
\end{annexitem}

\begin{annexitem}{annex:p15.clt-renewal-interpolation}{thm:p15-clt-renewal-interpolation}{The deterministic-window inverse calculation}
\begin{proof}
For large \(t\), \(u=\lfloor t/\mu\rfloor>M_t\).
Forward and reversed Kolmogorov inequalities give
\[
 \mathbb P\left(\max_{|m|\leq M_t}
       |T_{u+m}-T_u-m\mu|>t^{2/5}\right)
 \leq 2\sigma^2M_t/t^{4/5}\longrightarrow0.
 \]
The centered CLT for \(T_u\), or its variance bound, shows
\(D_t=T_u-t=O_{\mathbb P}(\sqrt t)\); the deterministic difference
\(\mu u-t\) lies in \([-\mu,0]\).
Consequently the event
\[
 |D_t|+t^{2/5}\leq\mu(M_t-2)
 \]
has probability tending to one. Intersect it with the event
of the preceding maximal bound. At the two window endpoints
the approximation \(T_{u+m}=t+D_t+\mu m+e_m\),
\(|e_m|\leq t^{2/5}\), shows that the count lies strictly
inside the window. Both \(N(t)\) and \(N(t)+1\) therefore lie
there. The inequalities
\(T_{N(t)}\leq t<T_{N(t)+1}\) then imply
\[
 |N(t)-(u-D_t/\mu)|\leq1+t^{2/5}/\mu.
 \]
Because \(|u-t/\mu|<1\), replacing \(u-D_t/\mu\) by
\((t-D_t)/\mu\) increases the bound by at most one.
Division by \(\sqrt t\) proves the second claimed convergence.
Moreover
\[
 \frac{D_t}{\sigma\sqrt{t/\mu}}
 =\sqrt{\frac{u}{t/\mu}}\,
       \frac{T_u-\mu u}{\sigma\sqrt u}
       +\frac{\mu u-t}{\sigma\sqrt{t/\mu}}
 \Rightarrow\chi.
 \]
The inverse estimate gives
\(\mu N(t)-t=-D_t+o_{\mathbb P}(\sqrt t)\);
normal symmetry finishes this distinct source proof.
\end{proof}
\end{annexitem}

\begin{annexitem}{annex:p15.clt-lindeberg-feller}{thm:p15-clt-lindeberg-feller}{The full finite-row transform proof}
\begin{proof}
If \(v_n\to0\), independence and centering give
\(\mathbb ES_n^2=v_n\to0\), proving that case.
Now suppose \(v_n\to\sigma^2>0\) and the Lindeberg condition.
For every \(\varepsilon>0\),
\[
 \max_jv_{n,j}
 \leq\varepsilon^2+L_n(\varepsilon),
 \]
so the maximum tends to zero by first taking \(n\to\infty\)
and then \(\varepsilon\downarrow0\).
For fixed \(t\in\mathbb R\), put
\(z_{n,j}=\mathbb Ee^{itX_{n,j}}\) and
\(w_{n,j}=1-t^2v_{n,j}/2\).
The second-order Taylor remainder proved in the characteristic
chapter bounds
\[
 |z_{n,j}-w_{n,j}|
 \leq\mathbb E\min(|tX_{n,j}|^3,2|tX_{n,j}|^2)
 \leq\varepsilon|t|^3v_{n,j}
       +2t^2\mathbb E(X_{n,j}^2;|X_{n,j}|>\varepsilon).
 \]
The absolute value on the cubic power is essential for negative \(t\).
Summing over the entire finite row yields
\[
 \limsup_n\sum_j|z_{n,j}-w_{n,j}|
 \leq\varepsilon|t|^3\sigma^2.
 \]
Let \(\varepsilon\downarrow0\).
Each \(|z_{n,j}|\leq1\), and eventually \(0<w_{n,j}\leq1\).
The product-comparison lemma with \(\theta=1\) therefore gives
\(\prod_jz_{n,j}-\prod_jw_{n,j}\to0\).

To justify the final real product limit explicitly, for \(0\leq u\leq1/2\)
the fundamental theorem of calculus gives
\[
 |\log(1-u)+u|
 =\int_0^u\frac{s}{1-s}\,ds\leq u^2.
 \]
With \(u_{n,j}=t^2v_{n,j}/2\), their maximum tends to zero,
\(\sum_ju_{n,j}\to t^2\sigma^2/2\), and
\(\sum_ju_{n,j}^2\leq(\max_ju_{n,j})\sum_ju_{n,j}\to0\).
Thus \(\prod_jw_{n,j}\to e^{-t^2\sigma^2/2}\).
Independence identifies \(\prod_jz_{n,j}\) with the transform
of \(S_n\), and the continuity theorem proves the limit.
All sums were over \(k_n\), so no bound on the growth of row length
and no relation between distinct rows was used.

For the source's i.i.d. embedding \(X_{n,j}=X_j/\sqrt n\),
the total variance is \(\sigma^2\), and
\[
 L_n(\varepsilon)
 =\mathbb E(X_1^2;|X_1|>\varepsilon\sqrt n)\to0
 \]
by dominated convergence. This explicitly verifies the claimed
inclusion of the earlier i.i.d. theorem.
\end{proof}
\end{annexitem}

\begin{annexitem}{annex:p15.clt-diagonal-truncation}{thm:p15-clt-diagonal-truncation}{Choosing the cutoff and estimating the removed sum}
\begin{proof}
Choose strictly increasing integers \(N_r\), tending to infinity,
such that \(L_n(1/r)\leq1/r\) for all \(n\geq N_r\).
This is possible by the Lindeberg hypothesis for each fixed \(r\),
and increasing a chosen threshold preserves its property.
Set \(\varepsilon_n=1\) before \(N_1\), and
\(\varepsilon_n=1/r\) on \(N_r\leq n<N_{r+1}\).
This is nonincreasing, tends to zero, and has
\(L_n(\varepsilon_n)\leq1/r\) on that block.

Put \(R_{n,j}=X_{n,j}\mathbf1_{\{|X_{n,j}|>\varepsilon_n\}}\).
Because \(\mathbb EX_{n,j}=0\),
\(\mathbb EY_{n,j}=-\mathbb ER_{n,j}\), and hence
\[
 S_n-U_n=\sum_j(R_{n,j}-\mathbb ER_{n,j}).
 \]
These centered remainders are independent within each row, so
\[
 \mathbb E(S_n-U_n)^2
 =\sum_j\operatorname{Var}(R_{n,j})
 \leq\sum_j\mathbb ER_{n,j}^2=L_n(\varepsilon_n)\to0.
 \]
Also \(|Y_{n,j}|\leq\varepsilon_n\) and
\(|\mathbb EY_{n,j}|\leq\varepsilon_n\), giving
\(|Z_{n,j}|\leq2\varepsilon_n\).
The \(L^2\) triangle inequality, itself a consequence of
Cauchy--Schwarz applied to the cross term, gives
\[
 \left|(\mathbb EU_n^2)^{1/2}-(\mathbb ES_n^2)^{1/2}\right|
 \leq(\mathbb E(S_n-U_n)^2)^{1/2}\to0.
 \]
Since both sums are centered, their second moments are their
total row variances. This proves the remaining assertion.
\end{proof}
\end{annexitem}

\begin{annexitem}{annex:p15.clt-moment-array-proof}{thm:p15-clt-moment-array-proof}{Complete moment combinatorics}
\begin{proof}
Write \(q_{n,j}=\mathbb EZ_{n,j}^2\) and
\(q_n=\sum_jq_{n,j}\to\sigma^2\).
For a fixed positive integer \(r\), expansion of
\((\sum_jZ_{n,j})^r\) groups the terms by their \(k\)
distinct indices and positive multiplicities
\(r_1+\cdots+r_k=r\):
\[
 \mathbb EU_n^r
 =\sum_{k=1}^r\ 
   \sum_{\substack{r_1,\ldots,r_k\geq1\\r_1+\cdots+r_k=r}}
 \frac{r!}{r_1!\cdots r_k!\,k!}
 \sum_{\substack{i_1,\ldots,i_k\leq k_n\\\text{all distinct}}}
      \prod_{\ell=1}^k\mathbb EZ_{n,i_\ell}^{r_\ell}.
 \]
The multinomial coefficient assigns positions among \(r\) factors;
the \(k!\) divides out the permutations of the ordered list of
distinct indices together with their multiplicities.
For \(k>k_n\) the inner sum is empty.
All products factor by row independence.
Any term with some \(r_\ell=1\) vanishes by centering.

When all \(r_\ell=2\), the inner sum is bounded above by
\(q_n^k\). The difference consists of tuples with at least
one repeated pair of indices. For each pair of positions,
\[
 \sum_{\text{tuples with }i_a=i_b}
             \prod_{\ell=1}^kq_{n,i_\ell}
 =\left(\sum_jq_{n,j}^2\right)q_n^{k-2}
 \leq b_n^2q_n^{k-1}\quad(k\geq2),
 \]
since \(q_{n,j}\leq b_n^2\). Summing over the
\(\binom{k}{2}\) pairs shows that the difference tends to zero.
For \(k=1\) it is exactly zero. Thus the distinct-index sum
tends to \(\sigma^{2k}\).

If all multiplicities are at least two and some exceeds two,
\[
 \left|\mathbb EZ_{n,j}^{r_\ell}\right|
 \leq b_n^{r_\ell-2}q_{n,j}
 \]
and the absolute inner sum is at most
\(b_n^{r-2k}q_n^k\to0\).
There are only finitely many multiplicity patterns for fixed \(r\).
For odd \(r\) none consists entirely of twos, so the limit is zero.
For \(r=2k\) the only surviving pattern has coefficient
\((2k)!/(2^kk!)\) and gives the stated normal moment.

For \(\sigma>0\), the normal moments determine the normal law by
the proved moment-determinacy theorem, and the method of moments
(Theorem~\ref{thm:p15-criteria-moment-method}) gives
\(U_n\Rightarrow\sigma\chi\).
For \(\sigma=0\), its second moment tends to zero directly.
Apply this result with \(b_n=2\varepsilon_n\) from the
diagonal-truncation lemma. The difference \(S_n-U_n\to0\) in
probability then transfers the normal limit to \(S_n\).
This argument establishes the source's second Lindeberg proof
without using the first proof to obtain moment convergence.
\end{proof}
\end{annexitem}

\begin{annexitem}{annex:p15.clt-records-and-cycles}{thm:p15-clt-records-and-cycles}{Bernoulli rows and the exact record ranks}
\begin{proof}
The independent Bernoulli sum has expectation \(H_n\) and variance
\(H_n-\sum_{j=1}^nj^{-2}\).
Integral comparison gives \(\log n\leq H_n\leq1+\log n\)
and \(\sum_{j=1}^\infty j^{-2}<\infty\).
The row entries
\((B_j-1/j)/\sqrt{\log n}\) have total variance tending to one
and absolute values at most \(1/\sqrt{\log n}\).
For every fixed positive cutoff, the Lindeberg tail sum is
therefore identically zero for all sufficiently large \(n\).
The array theorem proves the centered normal limit, and
\((H_n-\log n)/\sqrt{\log n}\to0\) gives the second version.
The cycle-count law is precisely this sum law by the already proved
Theorem~\ref{thm:p15-permutation-cycle-law}.

For the records, let \(V_1,V_2,\ldots\) be i.i.d. with continuous
CDF. Their law has no atoms. Product integration gives
\(\mathbb P(V_i=V_j)=\int\mathbb P(V_i=x)\,d\mathbb P_{V_j}(x)=0\)
for \(i\ne j\). Thus for each fixed \(n\) their ranks form
a permutation almost surely. Give rank one to the largest value.
All \(n!\) strict ordering events have the same probability:
permuting coordinates preserves the product law and bijects the
ordering events. They are disjoint and exhaustive outside the
null tie event, so each has probability \(1/n!\).

For \(m<n\), condition on any specified distinct ranks of
positions \(m+1,\ldots,n\). Among the \(m\) available ranks
the first \(m\) positions have all \(m!\) orders equally likely.
Exactly \((m-1)!\) put the best available rank at position \(m\).
The conditional chance of a record at \(m\) is therefore \(1/m\).
Every later record event is determined by those later ranks:
whether position \(j>m\) beats all its predecessors depends on
its rank and on which ranks occupy positions after \(j\).
Thus conditioning successively on later ranks gives, for
\(m_1<\cdots<m_k\),
\[
 \mathbb P(A_{m_1}\cap\cdots\cap A_{m_k})
       =\prod_{\ell=1}^k1/m_\ell,
 \quad A_j=\{V_j>\max_{i<j}V_i\},\quad A_1=\Omega.
 \]
The identities for all intersections imply independence of the
indicators: mixed intersections with complements follow by finite
inclusion-exclusion. This proves the complete record premise and
then the same normal limit.
\end{proof}
\end{annexitem}

\begin{annexitem}{annex:p15.clt-three-series-source}{thm:p15-clt-three-series-source}{The formerly deferred central-limit contradiction}
\begin{proof}
Almost-sure convergence implies \(X_j\to0\) there.
If \(\sum_j\mathbb P(|X_j|>A)=\infty\), the independent
Borel--Cantelli theorem gives \(|X_j|>A\) infinitely often with
probability one, a contradiction. Thus the first series is finite.
The first Borel--Cantelli theorem then gives \(X_j=Y_j\)
eventually almost surely, so \(\sum_jY_j\) converges almost surely.

Suppose \(c_n=\sum_{j=1}^n\operatorname{Var}(Y_j)\to\infty\).
For all sufficiently large \(n\), it is positive, and the centered
row \((Y_j-\mathbb EY_j)/\sqrt{c_n}\) has variance sum one
and absolute bound \(2A/\sqrt{c_n}\to0\).
The Lindeberg theorem gives
\[
 U_n=\frac{\sum_{j=1}^n(Y_j-\mathbb EY_j)}{\sqrt{c_n}}
                                               \Rightarrow\chi.
 \]
On the other hand
\(V_n=\sum_{j=1}^nY_j/\sqrt{c_n}\to0\) almost surely,
because its numerator converges to a finite limit.
The perturbation lemma would make the deterministic constants
\[
 U_n-V_n=-\frac{\sum_{j=1}^n\mathbb EY_j}{\sqrt{c_n}}
 \]
converge weakly to \(\chi\), which is impossible.
Therefore \(\sum_j\operatorname{Var}(Y_j)<\infty\).
The variance-series convergence theorem gives almost-sure
convergence of \(\sum_j(Y_j-\mathbb EY_j)\). Subtract this
convergent series from \(\sum_jY_j\) on their common event
of probability one. The resulting deterministic series
\(\sum_j\mathbb EY_j\) converges. This supplies the entire
source argument, including all cases of initial zero \(c_n\).
\end{proof}
\end{annexitem}

\begin{annexitem}{annex:p15.clt-infinite-variance-example}{thm:p15-clt-infinite-variance-example}{The logarithmic variance and the truncation calculation}
\begin{proof}
Write \(R=|X_1|\). Since \(\mathbb P(R>1)=1\),
\(\mathbb P(R>x)=1\) for \(0\leq x<1\).
The specified symmetric law exists: its Lebesgue density is
\(|x|^{-3}\mathbf1_{\{|x|>1\}}\), whose integral is
\(2\int_1^\infty x^{-3}\,dx=1\) and whose absolute tail is
exactly \(x^{-2}\) for \(x\geq1\).
The tail-integral identity gives
\[
 \mathbb ER=1+\int_1^\infty x^{-2}\,dx=2,\qquad
 \mathbb ER^2=1+\int_1^\infty2x^{-1}\,dx=\infty.
 \]
Symmetry and integrability imply \(\mathbb EX_1=0\).
For \(c\geq1\), integration of
\(\min(R^2,c^2)=\int_0^c2x\mathbf1_{\{R>x\}}\,dx\)
and subtraction of \(c^2\mathbf1_{\{R>c\}}\) give
\[
 \mathbb E(R^2;R\leq c)
 =\int_0^c2x\mathbb P(R>x)\,dx-c^2\mathbb P(R>c)
 =1+2\log c-1=2\log c.
 \]
The identity handles strict and closed truncation exactly; the
prescribed tail is continuous at every \(c\geq1\).

For the source's \(c_n=\sqrt n\log\log n\), take \(n\)
sufficiently large that \(c_n\geq1\).
The union bound gives \(n\mathbb P(R>c_n)=(\log\log n)^{-2}\).
For \(Y_{n,j}=X_j\mathbf1_{\{|X_j|\leq c_n\}}\), symmetry gives
mean zero and the exact second moment \(2\log c_n\sim\log n\).
The source's distinct estimates also hold: its second moment is at
most \(\int_0^{c_n}2x\mathbb P(R>x)\,dx=1+2\log c_n\).
For \(1\leq x\leq\sqrt n\),
\[
 \mathbb P(|Y_{n,1}|>x)
 =x^{-2}-c_n^{-2}
 \geq\{1-(\log\log n)^{-2}\}x^{-2}.
 \]
Integration on that interval gives the lower bound
\(\{1-(\log\log n)^{-2}\}\log n\).

The entries \(Y_{n,j}/\sqrt{n\log n}\) have total variance
\(2\log c_n/\log n\to1\), and absolute bound
\(\log\log n/\sqrt{\log n}\to0\).
Their Lindeberg tails eventually vanish for each positive cutoff.
The array theorem proves the normalized truncated-sum limit.
The probability that the original and truncated sums differ
tends to zero, so perturbation transfers the same limit.
\end{proof}
\end{annexitem}

\begin{annexitem}{annex:p15.clt-escaping-variance}{thm:p15-clt-escaping-variance}{Rare terms and the lost second moments}
\begin{proof}
On a product space let \(\varepsilon_j\) be independent fair signs,
and independently let \(B_j\) be independent Bernoulli variables
of parameter \(j^{-2}\). Put
\[
 X_j=\varepsilon_j\{1+(j-1)B_j\}.
 \]
This includes \(j=1\), where it is a fair sign, and gives the
specified laws. Since \(\sum_j\mathbb P(B_j=1)<\infty\),
only finitely many exceptional terms occur almost surely.
Therefore
\[
 \frac1{\sqrt n}\sum_{j=1}^n(X_j-\varepsilon_j)\longrightarrow0
                                                   \quad\text{almost surely}.
 \]
The fair-sign i.i.d. CLT transfers to \(S_n/\sqrt n\).
For every \(j\geq1\),
\(\mathbb EX_j^2=1-j^{-2}+j^2j^{-2}=2-j^{-2}\).
Independence gives
\(\operatorname{Var}(S_n)=2n-\sum_{j=1}^nj^{-2}\).

For fixed \(\varepsilon>0\) and sufficiently large \(n\),
the ordinary values \(\pm1\) are below \(\varepsilon\sqrt n\).
The tail contribution of an exceptional \(j>\varepsilon\sqrt n\)
to the normalized second moment is
\(n^{-1}j^2j^{-2}=1/n\).
Consequently the full Lindeberg tail sum is
\[
 \frac{\#\{j\leq n:j>\varepsilon\sqrt n\}}n\longrightarrow1.
 \]
This identifies the missing hypothesis directly.
\end{proof}
\end{annexitem}

\begin{annexitem}{annex:p15.clt-lyapunov}{thm:p15-clt-lyapunov}{The normalized moment bound}
\begin{proof}
Discard finitely many rows with zero variance if necessary, as
allowed by the eventual positivity assumption, and set
\(X_{n,j}=(W_{n,j}-\mathbb EW_{n,j})/a_n\).
These entries are centered, row-independent, with variance sum one.
For every \(\varepsilon>0\), the pointwise inequality
\[
 x^2\mathbf1_{\{|x|>\varepsilon\}}
       \leq\varepsilon^{-\delta}|x|^{2+\delta}
 \]
gives
\[
 L_n(\varepsilon)
 \leq\varepsilon^{-\delta}a_n^{-2-\delta}
       \sum_j\mathbb E|W_{n,j}-\mathbb EW_{n,j}|^{2+\delta}\to0.
 \]
The Lindeberg theorem proves the assertion for its full stated
row domain and every \(\delta>0\).

If \(|X_j|\leq M\) and cumulative variance \(a_n^2\to\infty\),
then \(|X_j-\mathbb EX_j|\leq2M\). Thus the centered normalized
entries are bounded by \(2M/a_n\to0\), so their Lindeberg tails
eventually vanish. Equivalently their Lyapunov sum is at most
\((2M/a_n)^\delta\). A case \(M=0\) cannot meet the divergent
variance hypothesis.
For the uniform moment bound with variances one,
\(a_n=\sqrt n\) and the Lyapunov expression is at most
\(Cn/n^{1+\delta/2}=Cn^{-\delta/2}\to0\).
Taking \(W_{n,j}=X_j\), \(k_n=n\), gives exactly the source's
independent-sequence Lyapunov formulation.
\end{proof}
\end{annexitem}

\begin{annexitem}{annex:p15.clt-sparse-jumps}{thm:p15-clt-sparse-jumps}{All three sparse-jump parameter regimes}
\begin{proof}
For \(\beta>1\),
\(\sum_j\mathbb P(X_j\ne0)=\sum_jj^{-\beta}<\infty\).
Borel--Cantelli proves eventual zero values and convergence.
For \(0<\beta<1\), all means are zero and
\(\operatorname{Var}(X_j)=j^{2-\beta}\). Integral comparison,
or Riemann sums for \(x^{2-\beta}\) on \([0,1]\), gives
\[
 n^{-(3-\beta)}\sum_{j=1}^nj^{2-\beta}
       =\frac1n\sum_{j=1}^n(j/n)^{2-\beta}
       \longrightarrow\frac1{3-\beta}.
 \]
The entries \(X_j/n^{(3-\beta)/2}\) have absolute values at
most \(n^{(\beta-1)/2}\to0\). The Lindeberg theorem gives
the stated normal limit with its exact variance.

Let \(\beta=1\). For fixed \(t\), independence gives
\[
 \mathbb Ee^{itS_n/n}
 =\prod_{j=1}^n(1-u_{n,j}),\qquad
 u_{n,j}=\frac{1-\cos(tj/n)}j.
 \]
Here \(0\leq u_{n,j}\leq t^2j/(2n^2)\), so
\(\max_ju_{n,j}\leq t^2/(2n)\to0\).
The function
\[
 f_t(x)=\frac{1-\cos(tx)}x\quad(x>0),\qquad f_t(0)=0,
 \]
is continuous on \([0,1]\), since
\(0\leq f_t(x)\leq t^2x/2\).
Thus
\(\sum_ju_{n,j}=n^{-1}\sum_j f_t(j/n)\to\int_0^1 f_t(x)\,dx\).
The logarithm remainder used in the array proof gives
\(\sum_j|\log(1-u_{n,j})+u_{n,j}|
\leq(\max_ju_{n,j})\sum_ju_{n,j}\to0\).
This proves convergence to the displayed characteristic function.
Its exponent lies between \(-t^2/4\) and zero, and tends to
zero at \(t=0\); hence the limit is continuous at zero.
The characteristic-function continuity theorem constructs a
proper probability law with this transform and proves weak convergence.

It is not normal. The exponential power series, integrated
termwise on \([0,1]\) using its uniform absolute convergence
on each bounded \(t\)-interval, gives
\[
 \log\mathbb Ee^{itQ}
       =-t^2/4+t^4/96+O(t^6)\quad(t\to0).
 \]
The positive quartic coefficient excludes a Gaussian transform,
whose logarithm near zero has degree at most two. It also
excludes a point mass. This proves the required borderline
law without invoking a later stable-law or infinitely divisible
distribution theorem.
\end{proof}
\end{annexitem}
